\subsubsection{Tensor products and multi-qubit probabilities}\label{ex:tensor-products-and-multi-qubit-probabilities}

\exercise (\emph{Expand a tensor product symbolically}) Given two states:
\begin{equation*}
    \ket{\psi_{1}} = \alpha_1 \ket{0} + \beta_1 \ket{1}, \quad
    \ket{\psi_{2}} = \alpha_2 \ket{0} + \beta_2 \ket{1}
\end{equation*}
Compute:
\begin{equation*}
    \ket{\Psi} = \ket{\psi_1} \otimes \ket{\psi_2}
\end{equation*}

\highspace
\solution The tensor product (\autopageref{subsubsec:tensor-product}) is a mathematical operation that takes two vectors and produces a new vector that represents the combined state of the two systems:
\begin{equation*}
    \begin{aligned}
        \ket{\Psi} &= \ket{\psi_1} \otimes \ket{\psi_2} \\
        &= \left(\alpha_1 \ket{0} + \beta_1 \ket{1}\right) \otimes \left(\alpha_2 \ket{0} + \beta_2 \ket{1}\right) \\
        &= \alpha_1 \alpha_2 \ket{00} + \alpha_1 \beta_2 \ket{01} + \beta_1 \alpha_2 \ket{10} + \beta_1 \beta_2 \ket{11}
    \end{aligned}
\end{equation*}

\highspace
\exercise (\emph{Expand a Concrete Two-Qubit State}) Now let's move from the symbolic formula to an actual computation. Given the two qubits:
\begin{equation*}
    \ket{\psi_1} = \frac{1}{\sqrt{2}}\left(\ket{0} + \ket{1}\right) \quad \text{and} \quad \ket{\psi_2} = \ket{1}
\end{equation*}
Compute the combined state:
\begin{equation*}
    \ket{\Psi} = \ket{\psi_1} \otimes \ket{\psi_2}
\end{equation*}
\solution As in the previous exercise, we will first substitute:
\begin{equation*}
    \begin{aligned}
        \ket{\Psi} &= \ket{\psi_1} \otimes \ket{\psi_2} \\
        &= \left(\frac{1}{\sqrt{2}}\left(\ket{0} + \ket{1}\right)\right) \otimes \ket{1} \\
        &= \frac{1}{\sqrt{2}}\left(\ket{0} \otimes \ket{1} + \ket{1} \otimes \ket{1}\right) \\
        &= \frac{1}{\sqrt{2}}\left(\ket{01} + \ket{11}\right)
    \end{aligned}
\end{equation*}

\highspace
\exercise (\emph{Full-state probabilities}) Given the state from the previous exercise:
\begin{equation*}
    \ket{\Psi} = \frac{1}{\sqrt{2}}\left(\ket{01} + \ket{11}\right)
\end{equation*}
Compute the probabilities:
\begin{equation*}
    P(q_1 q_2 = 00), \quad P(q_1 q_2 = 01), \quad P(q_1 q_2 = 10), \quad P(q_1 q_2 = 11)
\end{equation*}
\solution The probability of measuring a specific basis state is equal to the squared modulus of its amplitude in the quantum state. Since:
\begin{equation*}
    \ket{\Psi} = \frac{1}{\sqrt{2}}\left(\ket{01} + \ket{11}\right)
\end{equation*}
The amplitudes of the four computational basis states are:
\begin{center}
    \begin{tabular}{@{} c c @{}}
        \toprule
        State & Amplitude \\
        \midrule
        \ket{00} & 0 \\[.3em]
        \ket{01} & $\frac{1}{\sqrt{2}}$ \\[.3em]
        \ket{10} & 0 \\[.3em]
        \ket{11} & $\frac{1}{\sqrt{2}}$ \\
        \bottomrule
    \end{tabular}
\end{center}
Therefore:
\begin{equation*}
    \begin{aligned}
        P\left(q_1q_2=00\right) &= \left|0\right|^2 = 0 \\
        P\left(q_1q_2=01\right) &= \left|\frac{1}{\sqrt{2}}\right|^2=\frac{1}{2}=50\% \\
        P\left(q_1q_2=10\right) &= \left|0\right|^2 = 0 \\
        P\left(q_1q_2=11\right) &= \left|\frac{1}{\sqrt{2}}\right|^2=\frac{1}{2}=50\%
    \end{aligned}
\end{equation*}
As expected, the probabilities sum to one:
\begin{equation*}
    P(00)+P(01)+P(10)+P(11)
    =
    0+\frac{1}{2}+0+\frac{1}{2}
    =
    1
\end{equation*}

\highspace
\exercise (\emph{Compute the probability of measuring individual qubits}) We continue from:
\begin{equation*}
    \ket{\Psi} = \frac{1}{\sqrt{2}}\left(\ket{01} + \ket{11}\right)
\end{equation*}
From the previous exercise, we know the full-state probabilities:
\begin{equation*}
    P(00)=0, \quad P(01)=\frac{1}{2}, \quad P(10)=0, \quad P(11)=\frac{1}{2}
\end{equation*}
Now, compute the marginal probabilities for each qubit:
\begin{equation*}
    P(q_1=0), \quad P(q_1=1), \quad P(q_2=0), \quad P(q_2=1)
\end{equation*}
\solution The \definition{Marginal Probability} of a qubit being in a certain state is the \hl{sum of the probabilities of all the full states that include that qubit state.} In other words, it is the \textbf{probability of measuring a single qubit}, regardless of the value of the other qubit. Given a two-qubit system, the marginal probabilities are computed as follows:
\begin{equation*}
    \begin{aligned}
        P(q_1=0) &= \sum P(0-) = P(00) + P(01) = 0 + \frac{1}{2} = \frac{1}{2} \\
        P(q_1=1) &= \sum P(1-) = P(10) + P(11) = 0 + \frac{1}{2} = \frac{1}{2} \\
        P(q_2=0) &= \sum P(-0) = P(00) + P(10) = 0 + 0 = 0 \\
        P(q_2=1) &= \sum P(-1) = P(01) + P(11) = \frac{1}{2} + \frac{1}{2} = 1
    \end{aligned}
\end{equation*}
Where $P(0-)$ means the sum of probabilities of all states where the first qubit is 0, and the second qubit can be anything (denoted by the dash, \emph{don't care}). Similarly, $P(-0)$ means the sum of probabilities of all states where the second qubit is 0, and the first qubit can be anything. And so on for the other cases.